%HOMEWORK template for M 325K
\documentclass{article}
\headheight=7pt         \topmargin=14pt
\textheight=574pt       \textwidth=445pt
\oddsidemargin=18pt     \evensidemargin=18pt

%Begin package import

\usepackage{amsmath, amssymb, amsthm}
\usepackage{mathtools}
\usepackage{pdfpages}
\usepackage{tikz-cd}

%End package import
%Begin custom commands

\newcommand{\C}{\mathbb{C}}
\newcommand{\R}{\mathbb{R}}
\newcommand{\Q}{\mathbb{Q}}
\newcommand{\Z}{\mathbb{Z}}
\newcommand{\N}{\mathbb{N}}
\newcommand{\F}{\mathbb{F}}
\newcommand{\abs}[1]{\left\lvert #1\right\rvert}
\newcommand{\RM}[1]{\mathrm{#1}}
\newcommand{\BF}[1]{\textbf{#1}}
\newcommand{\e}{\epsilon}
\newcommand{\ceil}[1]{\lceil{#1}\rceil}
\newcommand{\floor}[1]{\lfloor{#1}\rfloor}
\newcommand{\rarr}[0]{\rightarrow}
\newcommand{\IM}[1]{\text{Im}(#1)}
\newcommand{\Hom}[0]{\text{Hom}}
\newcommand{\Pow}[1]{\text{Pow}(#1)}
\newcommand{\angles}[1]{\langle #1 \rangle}

%End custom commands
%Begin custom theorems

\newtheorem{innercustomgeneric}{\customgenericname}
\providecommand{\customgenericname}{}
\newcommand{\newcustomtheorem}[2]{%
\newenvironment{#1}[1]
{%
\renewcommand\customgenericname{#2}%
\renewcommand\theinnercustomgeneric{##1}%
\innercustomgeneric%
}
{\endinnercustomgeneric}
}

\newcustomtheorem{THRM}{Theorem}
\newcustomtheorem{LEM}{Lemma}
\newcustomtheorem{EX}{Exercise}
\newcustomtheorem{COR}{Corollary}
\newcustomtheorem{PROB}{Problem}
\newcustomtheorem{claim}{Claim}

\newenvironment{solution}
{\renewcommand\qedsymbol{$\blacksquare$}\begin{proof}[Solution]}
{\end{proof}}

\newenvironment{definition}
{\renewcommand\qedsymbol{$\blacksquare$}\begin{proof}[Definition]}
{\end{proof}}

%End custom theorems
\begin{document}

\title{Chapter 1}
\author{Arham Lodha}%put your name here
\date{May 08, 2024}%enter the due date here
\maketitle

\section*{1.6. Cardinality}
% \begin{claim}{1.6}\label{Claim 1.6.}
%     Let $2 \N$ denote the even positive integers. Then $\abs{2 \N} = \abs{\N}$  
% \end{claim}
% \begin{proof}
%     Let $f: \N \to 2 \N$ where $n \to 2n$. $\forall m \in 2 \N, m / 2 \in \N, f(m / 2) = m$. $\forall p, q \in \N, 2p = f(p) = f(q) = 2q \implies p = q $. Thus $f$ is injective and surjective. Thus $\abs{2\N} = \abs{\N}.$        
% \end{proof}
% \bigskip
% \begin{PROB}{1.7}\label{Problem 1.7.}
%     The set $\Z$ has the same cardinality as $\N$, that is, $|\Z| = |N|$.
% \end{PROB}
% \begin{proof}
%     Let $f: \N \to \Z$, be defined as the following, $\forall x \in \N$:
    
%     \begin{equation*}
%         f(x) = \begin{cases*}
%             \frac{x}{2} & \text{if x is even} \\ 
%             -\frac{x - 1}{2} & \text{otherwise}
%         \end{cases*}
%     \end{equation*}

%     $\forall n, m \in \N$, $f(n) = f(m)$. If $f(n) \leq 0 \implies f(m) \leq 0$, thus $-\frac{n - 1}{2} = -\frac{m - 1}{2} \implies n = m$. Otherwise if $f(n), f(m) > 0 \implies \frac{n}{2} = \frac{m}{2} \implies n =m$. Thus $f$ is injective. 
    
%     $\forall z \in \Z$, if $z \leq 0$. Let $n = -2z + 1$ and $n \in \N$ and $f(n) = z$. If $z > 0$, let $n = 2z$ and thus $f(n) = z$. Thus $f$ is surjective. 
    
%     Thus f is a bijection and $\abs{\N} = \abs{\Z}$.  
% \end{proof}

% \bigskip

\begin{PROB}{1}\label{Problem 1.}
    Show $\Q$ is countably infinite 
\end{PROB}
\begin{proof}
    It suffices to show there is a injective map from $\Q \to \N \times \N$ since $\N \times \N$ is countable and $\Q$ is infinite. Define $f: \Q \to \N \times \N $ where $\forall \frac{p}{q} \in \Q \ni p, q \in \Z \ni q \neq 0$ where wlog $gcd(p, q) = 1$. 
    
    \begin{equation*}
        f\left(\frac{p}{q}\right) := \begin{cases*}
            (2p, 2q) & \text{if $p, q \geq 0$} \\
            (2p, -2q + 1) & \text{if $p\geq 0$ and $q < 0$ } \\
            (-2p + 1, 2q) & \text{if $q \geq 0$ and $p < 0$ } \\
            (-2p + 1, -2q + 1) & \text{if $p, q < 0$ } \\
        \end{cases*}
    \end{equation*}

    $\forall \frac{p}{q}, \frac{r}{s} \in \Q \ni f(\frac{p}{q}) = f(\frac{r}{s})$. 
    
    \begin{enumerate}
        \item If $p, q \geq 0 \implies r, s \geq 0$, thus $2p = 2r \implies p = r$, and $2q = 2s \implies q = s$. 
        \item Cases 2, 3, and 4 are the same.    
    \end{enumerate}
    
    f is injective. Thus $\Q$ is countably infinite. 
\end{proof}

\bigskip

\begin{PROB}{2}\label{Problem 2.}
    Show that the maps f and g of Examples 1 and 2 are bijections.
\end{PROB}
\begin{claim}{2a}\label{Claim 2a.}
    $f$ is a bijection. 

    \begin{equation*}
        f(n) := \begin{cases*}
            2n & \text{if $n > 0$} \\
            -2n + 1 & \text{if $n \leq 0$}
        \end{cases*}
    \end{equation*} 
\end{claim}
\begin{proof}
    $\forall n, m \in \Z \ni f(n) = f(m)$, if $n > 0 \implies 2n = f(m)$. Assume the contrary that $m \leq 0$, then $2n = -2m + 1 \implies m = n - \frac{1}{2} \implies m \not \in \Z$ which is a contradiction. Thus $m > 0$. Thus $2n = 2m \implies m = n$. If $n \leq 0 \implies -2n + 1 = f(m)$. Assume the contrary that $m > 0 \implies 2n - 1 =2m \implies m = n - \frac{1}{2} \implies m \not \in \Z$ which is a contradiction. Thus $2n - 1 = 2m - 1 \implies n = m$. Thus $f$ is a injective function. 
    
    $\forall y \in \Z_+$. 
    
    \begin{enumerate}
        \item If $y$ is even: Let $x = \frac{y}{2}$. $f(x) = y$.
        \item If $y$ is odd: Let $x = -\frac{y - 1}{2}$, $f(x) = -2x + 1 = y$.     
    \end{enumerate}

    Thus $f$ is surjective. 
\end{proof}
% \begin{claim}{2b}\label{Claim 2b.}
%     $g$ is a bijection.
    
%     \begin{equation*}
%         g(x, y) = \frac{1}{2}(x - 1)x + y
%     \end{equation*}
% \end{claim}
% \begin{proof}
%     $\frac{1}{2}(x - 1)(x) = \sum_{k = 1}^{x - 1} k = T_{x - 1}$. 

%     $\forall (x, y), (a, b) \in \Z_+ \times \Z_+$. $\frac{1}{2}(x - 1)x + y = g(x, y) = g(a, b) = \frac{1}{2}(a - 1)a + b \implies T_{x - 1} + 2y = T_{a - 1} + 2b$
% \end{proof}

\bigskip

\begin{PROB}{3}\label{Problem 3.}
    Let X be the two-element set $\left\{ 0, 1 \right\}$. Show there is a bijective correspondence between the set $2^{\Z_+}$ and the cartesian product $X^\omega$.
\end{PROB}
\begin{proof}
    Let $f : 2^{\Z_+} \to X^\omega$ where $A \in 2^{\Z_+}$ gets sent to a element of $X^\omega$  where the $i$th component is 1 iff $i \in A$ otherwise the ith component is 0.
    
    $\forall A, B \in 2^{\Z_+}, f(A) = f(B)$. $\forall i \in A$, then the ith component of $f(A)$ is 1. Thus the ith component of $f(B)$ is 1. Thus $i \in B$. Thus $A \subseteq B$. $\forall j \in B$, then $j$th component of $f(B)$ is 1. Thus the jth component of $f(A)$ is 1. THus $j \in A$. Thus $B \subseteq A$. $A = B$. Thus $f$ is injective. 
    
    $\forall (a_1, \dots, a_i, \dots) \in X^{\omega}$. Let $A = \left\{ x \in \Z_+ \mid a_x = 1 \right\}$. $f(A) = (a_1, \dots, a_i, \dots)$. Thus $f$ is surjective. 
    
    Thus $f$ is a bijection. 
\end{proof}

\section*{Infinite Sets and the Axiom of Choice}
\begin{PROB}{9.1}\label{Problem 9.1.}
    Define an injective map $f : \Z_+ \to X^\omega$, where X is the two-element set ${0,1}$ without using the choice axiom.
\end{PROB}
\begin{proof}
    Let $f := \Z_+ \to X^\omega$ where $m \to (\delta(1 - m), ..., \delta(i - m), \dots)$. 
    
    Suppose $\forall m, n \in \Z_+ \ni f(m) = f(n)$ that means $f(m)$ and $f(n)$ agree on the $m$ and $n$th entry. That means there are 1s on the m and nth entry. For $f(m)$, there is a 1 only on the mth entry similarly for $f(n)$ will only have a 1 on the nth entry. Thus there is a 1 on the mth entry and nth entry which means $m = n$.        
\end{proof}

\bigskip

\begin{PROB}{9.2}\label{Problem 9.2.}
    Find if possible a choice function for each of the following collections, without using the choice axiom:
    \begin{enumerate}
        \item The collection $A$ of nonempty subsets of $\Z_+$
        \item The collection $B$ of nonempty subsets of $\Z$
        \item The collection $C$ of nonempty subsets of the rational $\Q$
        \item The collection $D$ of nonempty subsets of $X^\omega$ where $X = \left\{ 0, 1 \right\}$
    \end{enumerate}
\end{PROB}
\begin{proof}
    \begin{enumerate}
        \item Let $c : A \to \bigcup_{B \in A} B \ni B \to \min(B)$.
        \item Let $f$ be the bijection from $\Z \to \Z_+$ seen in problem 1.6.1. Let $c : B \to \bigcup_{A \in B} A \ni A \to \min(f(A))$      
    \end{enumerate}

    For the rest do the same. First map to $\Z_+$ and then take the min which works since $\Z_+$ is well ordered.  
\end{proof}

\bigskip

\begin{PROB}{9.3}\label{Problem 9.3.}
    Suppose that A is a set and $\left\{ f_n \right\}_{n \in \Z_+}$ is a given indexed family of injective functions

    \[
        f_n := [1, \dots, n] \to A
    \]
    
    Show $A$ is infinite. Can you define a injective function from $f \to \Z_+ \to A$ without using the choice axiom.   
\end{PROB}
\begin{proof}
    Assume the contrary A is finite. Let $m := \abs{A}$. There is a bijective map $f_{m}$ from $ [1, \dots, m]\to A$. Thus since $f_{m + 1}$ is injective. There is a injective map from $[1, \dots, m+1] \to [1, \dots, m]$, which is a contradiction. Thus $A$ is infinite.
    
    Define $f$ recursively. Let $a_0 = f_1(1)$ and let 
    \[\rho(h : [1, \dots, m] \to A) = f_{m}(\min{k \in \Z_+ | k \leq n + 1, f_{m}(k) \not \in h([1, \dots, m - 1])})\] 

    Since $\rho$ is well defined $f_{n}$ is injective. Thus $f$ is defined in the recursive function and thus injective.    
\end{proof}

\bigskip

\section*{Well Order Sets}


\end{document}